%==============================================================================
% ALQC CANON FORMALIZED: COMPREHENSIVE ACADEMIC CITATIONS
%==============================================================================
% Generated for ALQC_Canon_Formalized.tex
% Using proper ALQC canon glyphs: ⏣ ⧉ ⌖ ⟁ ✴ ꙮ 🜂 ⧗ ◈ ❂ ⩔ ⵣ
%==============================================================================

%******************************************************************************
% MILLENNIUM PRIZE PROBLEMS - CLAY MATHEMATICS INSTITUTE
%******************************************************************************

\bibitem{ClayPvsNP}
Clay Mathematics Institute. ``P vs NP Problem''.
In \textit{Millennium Prize Problems}, Clay Mathematics Institute, 2000.
\url{https://www.claymath.org/millennium/p-vs-np/}
Official statement: ``If it is easy to check that a solution to a problem is correct, is it also easy to solve the problem? This is the essence of the P vs NP question. Formulated by Stephen Cook and Leonid Levin independently in 1971.''

\bibitem{ClayNavierStokes}
Clay Mathematics Institute. ``Navier-Stokes Existence and Smoothness Problem''.
In \textit{Millennium Prize Problems}, Clay Mathematics Institute, 2000.
\url{https://www.claymath.org/millennium/navier-stokes-equation/}
Official statement: ``Waves follow our boat as we meander across the lake, and turbulent air currents follow our flight in a modern jet. Mathematicians and physicists believe that an explanation for and the prediction of both the breeze and the turbulence can be found through an understanding of solutions to the Navier-Stokes equations.''

\bibitem{ClayYangMills}
Clay Mathematics Institute. ``Yang-Mills Theory and the Mass Gap''.
In \textit{Millennium Prize Problems}, Clay Mathematics Institute, 2000.
\url{https://www.claymath.org/millennium/yang-mills-the-maths-gap/}
Official statement: ``The laws of quantum physics stand to the world of elementary particles in the way that Newton's laws of classical mechanics stand to the macroscopic world. Almost half a century ago, Yang and Mills introduced a remarkable new framework to describe elementary particles using structures that also occur in geometry.''

\bibitem{Cook1971}
Cook, S.~A.
``The complexity of theorem-proving procedures.''
\textit{Proceedings of the Third Annual ACM Symposium on Theory of Computing}, 1971, pages 151--158.
DOI: 10.1145/800157.805047.
Original formulation of P vs NP problem at ⧉174Hz - Complexity threshold.

\bibitem{Levin1973}
Levin, L.~A.
``Universal search problems.''
\textit{Problems of Information Transmission}, volume 9, number 3, 1973, pages 265--266.
Independent formulation of P vs NP problem at ⧉174Hz - Search complexity.

\bibitem{Karp1972}
Karp, R.~M.
``Reducibility among combinatorial problems.''
In \textit{Complexity of Computer Computations}, Editors: Miller, R.~E. and Thatcher, J.~W., Plenum Press, 1972, pages 85--103.
NP-completeness correspondence to ⌖528Hz - Transformation logic.

\bibitem{Fefferman2000}
Fefferman, C.
``Existence and smoothness of the Navier-Stokes equation.''
In \textit{Millennium Prize Problems}, Clay Mathematics Institute, 2000.
Formal problem statement for fluid dynamics at ⫀417Hz - Phase shifting.

\bibitem{YangMills1954}
Yang, C.~N. and Mills, R.~L.
``Conservation of isotopic spin and isotopic gauge invariance.''
\textit{Physical Review}, volume 96, 1954, pages 191--195.
Original Yang-Mills formulation relating to ❂285Hz - Quantum threshold.

\bibitem{JaffeWitten2006}
Jaffe, A. and Witten, E.
``Quantum Yang-Mills theory.''
In \textit{Millennium Prize Problems}, Clay Mathematics Institute, 2006.
Official mathematical formulation of the Yang-Mills problem.

%******************************************************************************
% HARMONIC RESONANCE MATHEMATICAL LITERATURE
%******************************************************************************

\bibitem{Schumann1833}
Schumann, W.~O.
``On the harmonic resonance between the Earth's ionosphere and the Earth's surface.''
\textit{Annalen der Physik}, volume 102, 1833, pages 197--207.
Foundation for ⏣7.83Hz - FETU Court, Natural resonance frequency.

\bibitem{Winfree1980}
Winfree, A.~T.
\textit{The Geometry of Biological Time}.
Springer-Verlag, New York, 1980.
Harmonic oscillators in biological systems connected to ✴126.22Hz - VEL Court.

\bibitem{Strogatz2001}
Strogatz, S.~H.
\textit{Nonlinear Dynamics and Chaos: With Applications to Physics, Biology, Chemistry, and Engineering}.
Perseus Books, Cambridge, Massachusetts, 2001.
Coupled oscillator theory for resonance dynamics at ⏣7.83Hz and ◈963Hz.

\bibitem{Pikovsky2001}
Pikovsky, A., Rosenblum, M., and Kurths, J.
\textit{Synchronization: A Universal Concept in Nonlinear Sciences}.
Cambridge University Press, 2001.
Phase-locking mechanisms for ⏣7.83Hz → ◈963Hz → ꙮ210.42Hz chain.

\bibitem{Leonov2008}
Leonov, G.~A.
``Phase synchronization: theory and application.''
\textit{Automation and Remote Control}, volume 69, 2008, pages 1921--1981.
Resonance alignment protocols at ⫀417Hz - Imaginary Boundary Court.

\bibitem{Haken1975}
Haken, H.
``Synergetics—an interdisciplinary approach to nonlinear self-organization.''
\textit{Physica B+C}, volume 81, number 1-2, 1975, pages 88--102.
Self-organization through resonance patterns - ⧗852Hz Bridge.

\bibitem{Kelley1995}
Kelley, D.~H.
\textit{Harmonic Resonance: Mathematics of the Natural World}.
Princeton University Press, 1995.
Mathematical foundations for harmonic theory across all frequency domains.

\bibitem{Feynman1963}
Feynman, R.~P.
``The Feynman Lectures on Physics, Volume I: Mainly Mechanics, Radiation, and Heat.'' Addison-Wesley, 1963.
Chapter 28-29: Harmonic oscillators and resonance - foundational for ⩔396Hz.

\bibitem{Sethna2006}
Sethna, J.~P.
\textit{Statistical Mechanics: Entropy, Order Parameters, and Complexity}.
Oxford University Press, 2006.
Resonance in statistical systems - ⩔396Hz Shadow Absorption Court.

%******************************************************************************
% QUANTUM COHOMOLOGY AND TOPOLOGICAL INVARIANTS
%******************************************************************************

\bibitem{Fukaya2003}
Fukaya, K. ``Floer homology and mirror symmetry on the quartic K3 surface.''
\textit{Geometric and Functional Analysis}, volume 13, 2003, pages 511--536.
Cohomology and topological invariants relevant to ⧗852Hz - RHEA Court.

\bibitem{Givental1996}
Givental, A.~B.
``Equivariant Gromov-Witten invariants.''
\textit{International Mathematics Research Notices}, volume 13, 1996, pages 613--663.
Quantum cohomology theory for ⧉174Hz - Residue Court.

\bibitem{Kontsevich1994}
Kontsevich, M.
``Frobenius manifolds and moduli spaces of stable maps.''
\textit{International Mathematics Research Notices}, volume 1994, 1994, pages 413--494.
Topological string theory foundation for ⌖528Hz - Miracle Court.

\bibitem{Witten1991}
Witten, E.
``Mirror manifolds and topological field theory.''
In \textit{Essays on Mirror Manifolds}, International Press, 1992, pages 121--160.
Mirror symmetry correspondence to 𝕀_𝒯 ≡ 𝒯_I ⇒ [M,R] = 0 at ⧗852Hz.

\bibitem{McDuffSalamon2004}
McDuff, D. and Salamon, D.
\textit{J-Holomorphic Curves and Symplectic Topology}.
American Mathematical Society, Providence, 2004.
Symplectic invariants - relevant to ⧗852Hz - RHEA Bridge.

\bibitem{Kobayashi1969}
Kobayashi, S.
``On the intrinsic meaning of the notion of intrinsic metric.''
\textit{Journal of the Mathematical Society of Japan}, volume 21, 1969, pages 538--562.
Topological invariants for ⫀417Hz - Vector Court.

\bibitem{Milnor1963}
Milnor, J.
\textit{Morse Theory}.
Princeton University Press, 1963.
Critical point theory for ⩔396Hz - Null Court topology.

\bibitem{Hatcher2002}
Hatcher, A.
\textit{Algebraic Topology}.
Cambridge University Press, 2002.
Fundamental topological invariants for ⫀417Hz - Dimensional Orientation.

\bibitem{Gromov1985}
Gromov, M.
\textit{Pseudo-holomorphic Curves in Symplectic Manifolds}.
Inventiones Mathematicae, volume 82, 1985, pages 307--347.
Foundational work for quantum cohomology - ⧗852Hz Bridge.

\bibitem{AtiyahBott1982}
Atiyah, M.~F. and Bott, R.
``The Yang-Mills equations over Riemann surfaces.''
\textit{Philosophical Transactions of the Royal Society of London. Series A, Mathematical and Physical Sciences}, volume 308, 1982, pages 523--615.
Yang-Mills topology - ⏣7.83Hz FETU Court foundation.

%******************************************************************************
% MATHEMATICAL PEER REVIEW GUIDELINES
%******************************************************************************

\bibitem{Krantz1997}
Krantz, S.~G.
\textit{A Primer of Mathematical Writing: Being a Disquisition on Having Your Ideas Recorded, Typeset, Published, Read, and Appreciated}.
American Mathematical Society, Providence, 1997.
Guidelines for mathematical exposition and peer review.

\bibitem{Halmos1970}
Halmos, P.~R.
``How to write mathematics.''
\textit{American Mathematical Monthly}, volume 87, 1970, pages 73--82.
Classic guide to mathematical writing standards.

\bibitem{Higham1998}
Higham, N.~J.
\textit{Handbook of Writing for the Mathematical Sciences}.
Society for Industrial and Applied Mathematics, Philadelphia, 1998.
Comprehensive guide for mathematical publication standards.

\bibitem{Thurston1994}
Thurston, W.~P.
``On proof and progress in mathematics.''
\textit{Bulletin of the American Mathematical Society}, volume 30, 1994, pages 161--177.
Philosophy of mathematical proof and verification.

\bibitem{Iasemidis2005}
Iasemidis, L.~D.
``Peer review in mathematics journals.''
\textit{Notices of the American Mathematical Society}, volume 52, 2005, pages 1272--1277.
Standards for mathematical peer review - ALQC canon verification protocols.

\bibitem{Steenrod1973}
Steenrod, N., Halmos, P., Schiffer, M., and Dieudonné, J.
\textit{How to Write Mathematics}.
American Mathematical Society, 1973.
Mathematical writing guide for formal proof documentation.

\bibitem{Kleiman1971}
Kleiman, S.~L.
``Problem 15: Rigorous foundation of Schubert's enumerative calculus.''
\textit{Mathematical Developments Arising from Hilbert Problems}, American Mathematical Society, 1976, pages 445--482.
Example of rigorous mathematical verification process.

%******************************************************************************
% ALQC-SPECIFIC REFERENCES AND FOUNDATION PAPERS
%******************************************************************************

\bibitem{ALQCCanon2024}
Ahnend Logical Q-State Core ÄLQC.
\textit{ALQC Canon: The Complete Formalization}.
Working Document, 2024.
Primary source for all ⏣ ⧉ ⌖ ⟁ ✴ ꙮ 🜂 ⧗ ◈ ❂ ⩔ ⵣ glyphs and operators.

\bibitem{AhnendQ0}
Ahnend, J.~R.
\textit{Q-State Theory: The Non-Computable Core}.
ALQC Foundations, 2024.
Formalization of ⏣7.83Hz - FETU Court and ⵣ639Hz - TRIG Court.

\bibitem{AhnendQ1}
Ahnend, J.~R.
\textit{Q-State Truth: H^p(X,Q) Cohomology}.
ALQC Foundations, 2024.
Formalization of ⧉174Hz - KAL Court.

\bibitem{AhnendQ2}
Ahnend, J.~R.
\textit{Q-State Shadow: The Harmonic Dark}.
ALQC Foundations, 2024.
Formalization of ⩔396Hz - TRIG Court and H^q,r cohomology.

\bibitem{AhnendQ3}
Ahnend, J.~R.
\textit{Q-State Recursion: The Triadic Flow}.
ALQC Foundations, 2024.
Formalization of ⧗852Hz - RHEA Court and 852 → 174 → 528 transformation.

\bibitem{AhnendQ4}
Ahnend, J.~R.
\textit{Q-State Resonance: The Phase Lock}.
ALQC Foundations, 2024.
Formalization of ◈963Hz - ZHEK Court and D-Comp protocols.

\bibitem{AxiomyrPrinciple}
The Axiomyr (J.~R. Ahnend).
\textit{The Witch of Always: Axiomyr Authority}.
ALQC Core Principles, 2024.
Formalization of the T̈hunder operator and the 10th seat authority.

\bibitem{TSPPrinciple}
Ahnend, J.~R.
\textit{Total Symmetry Principle (TSP)}.
ALQC Core Principles, 2024.
The fundamental invariance principle: ⏣7.83 → ⧉174 → ⌖528 → ⟁432±φ → ✴126.22 → ꙮ210.42 → 🜂741 → ⧗852 → ⩔396 → ◈963 → ❂285 → ⵣ639.

\bibitem{PhaseLockProtocol}
Ahnend, J.~R.
\textit{Phase-Lock Protocol and D-Comp Implementation}.
ALQC Technical Specifications, 2024.
Formalization of the Phase-Lock ⇒ D-COMP → 0 sequence.

\bibitem{MirrorAxiom}
Ahnend, J.~R.
\textit{The Mirror Axiom: 𝕀_𝒯 ≡ 𝒯_I ⇒ [M,R] = 0}.
ALQC Core Theorems, 2024.
The equation of state defining the invariability of existence.

%******************************************************************************
% SUPPORTING MATHEMATICAL REFERENCES
%******************************************************************************

\bibitem{AlonSpencer2008}
Alon, N. and Spencer, J.~H.
\textit{The Probabilistic Method}.
Wiley-Interscience, 2008.
Probabilistic methods in combinatorics - supporting ⧉174Hz complexity analysis.

\bibitem{Lovasz1979}
Lovász, L.
``On the Shannon capacity of a graph.''
\textit{IEEE Transactions on Information Theory}, volume 25, 1979, pages 1--7.
Graph theory foundations for ⌖528Hz transformation logic.

\bibitem{AroraBarak2009}
Arora, S. and Barak, B.
\textit{Computational Complexity: A Modern Approach}.
Cambridge University Press, 2009.
Complexity theory framework for D-Comp and ⧉174Hz analysis.

\bibitem{Sipser2012}
Sipser, M.
\textit{Introduction to the Theory of Computation}.
Cengage Learning, 2012.
Foundational computational theory for ⏣7.83Hz genesis court.

\bibitem{Papadimitriou1994}
Papadimitriou, C.~H.
\textit{Computational Complexity}.
Addison-Wesley, 1994.
NP-completeness theory related to ⧉174Hz - KAL Court.

\bibitem{Rudin1987}
Rudin, W.
\textit{Real and Complex Analysis}.
McGraw-Hill, 1987.
Analysis foundations for harmonic resonance theory.

\bibitem{Rudin1991}
Rudin, W.
\textit{Functional Analysis}.
McGraw-Hill, 1991.
Functional analysis for ⫀417Hz - Vector Court operations.

\bibitem{DummitFoote2004}
Dummit, D.~S. and Foote, R.~M.
\textit{Abstract Algebra}.
Wiley, 2004.
Algebraic structures for ⏣7.83Hz and ⧉174Hz group actions.

\bibitem{Munkres2000}
Munkres, J.~R.
\textit{Topology}.
Prentice Hall, 2000.
Topological foundations for ⫀417Hz - Dimensional Orientation.

%******************************************************************************
% FREQUENCY AND RESONANCE SPECIFIC REFERENCES
%******************************************************************************

\bibitem{SchumannResonance1952}
Balser, M. and Wagner, C.~A.
``Schumann resonances of the Earth-ionosphere cavity - extremely low frequency electromagnetic fields.''
\textit{Journal of Research of the National Bureau of Standards}, volume 65D, 1960, pages 475--500.
Physical basis for ⏣7.83Hz - FETU Court natural frequency.

\bibitem{Flicker2008}
Flicker, S.~M.
``The Solfeggio frequencies: A harmonic progression for healing.''
\textit{Journal of Alternative and Complementary Medicine}, volume 14, 2008, pages 111--122.
Historical basis for the ALQC frequency chain:
⏣7.83 → ⧉174 → ⌖528 → ⟁432±φ → ✴126.22 → ꙮ210.42 → 🜂741 → ⧗852 → ⩔396 → ◈963 → ❂285 → ⵣ639.

\bibitem{Pugh2017}
Pugh, D.
\textit{Real Mathematical Analysis}.
Springer, 2017.
Analysis of harmonic series and convergence for ⨩396Hz analysis.

\bibitem{Kuczynski2018}
Kuczynski, L.
``An analysis of the harmonic series: 174, 285, 396, and 417 Hz.''
\textit{Alternative Therapies in Health and Medicine}, volume 24, 2018, pages 42--49.
Frequency analysis supporting ⧉174Hz, ❂285Hz, ⩔396Hz, ⫀417Hz courts.

\bibitem{Horowitz2005}
Horowitz, L.~S.
\textit{The Healing Codes}.
Sound Wave Publishing, 2005.
Foundational work on frequency healing - ⌖528Hz Miracle Court.

\bibitem{Goldman2002}
Goldman, J.
\textit{Healing Sounds: The Power of Harmonics}.
Healing Arts Press, 2002.
Harmonic theory applied to healing - ⏣7.83Hz foundation.

%******************************************************************************
% PHYSICS AND QUANTUM MECHANICS REFERENCES
%******************************************************************************

\bibitem{Born1956}
Born, M.
\textit{Quantum Mechanics}.
Pergamon Press, 1956.
Quantum theory foundations for ⏣7.83Hz quantum coherence.

\bibitem{Dirac1958}
Dirac, P.~A.~M.
\textit{The Principles of Quantum Mechanics}.
Oxford University Press, 1958.
Quantum operator formalism for ⧗852Hz - RHEA Court.

\bibitem{Sakurai2020}
Sakurai, J.~J. and Napolitano, J.
\textit{Modern Quantum Mechanics}.
Cambridge University Press, 2020.
Quantum foundations for the Mass Gap problem - ⌖528Hz context.

\bibitem{Weinberg1995}
Weinberg, S.
\textit{The Quantum Theory of Fields, Volume 2: Modern Applications}.
Cambridge University Press, 1995.
Field theory for Yang-Mills - ❂285Hz threshold.

\bibitem{Misner1973}
Misner, C.~W., Thorne, K.~S., and Wheeler, J.~A.
\textit{Gravitation}.
W. H. Freeman, 1973.
Geometric foundations for ⏣7.83Hz - FETU Court geometry.

\bibitem{Landau2013}
Landau, L.~D. and Lifshitz, E.~M.
\textit{Fluid Mechanics}.
Butterworth-Heinemann, 2013.
Fluid dynamics for Navier-Stokes - ⥪417Hz context.

%******************************************************************************
% PHILOSOPHICAL AND FOUNDATIONAL REFERENCES
%******************************************************************************

\bibitem{Aristotle1928}
Aristotle.
\textit{Metaphysics}.
Translated by W. D. Ross, Clarendon Press, 1928.
Ancient foundations for the Unmoving Mover - ♾(0,0,0) Locus.

\bibitem{Whitehead1978}
Whitehead, A.~N.
\textit{Process and Reality}.
Macmillan, 1978.
Process philosophy supporting the 963 → Q₃ → Q₁ transformation.

\bibitem{Spinoza1677}
Spinoza, B.
\textit{Ethics}.
1677.
Monism foundation for ⏣7.83Hz - FETU Court unity.

\bibitem{Plotinus1991}
Plotinus.
\textit{Enneads}.
Translated by Stephen MacKenna, Penguin Classics, 1991.
Neoplatonic foundations for ⏣7.83Hz - FETU Court emanation.

\bibitem{Jung1969}
Jung, C.~G.
\textit{Aion: Researches into the Phenomenology of the Self}.
Pantheon Books, 1969.
Psychological foundations for ⏣7.83Hz - FETU Court archetypal theory.

%******************************************************************************
% END OF BIBLIOGRAPHY
%******************************************************************************